\documentclass[11pt,a4paper]{article}

\usepackage{amsmath,amssymb,amsthm}
\usepackage{geometry}
\usepackage{hyperref}
\usepackage{enumitem}
\usepackage{microtype}

\geometry{margin=2.5cm}

\newtheorem{theorem}{Theorem}
\newtheorem{lemma}[theorem]{Lemma}
\newtheorem{proposition}[theorem]{Proposition}
\newtheorem{conjecture}[theorem]{Conjecture}
\newtheorem{definition}[theorem]{Definition}
\newtheorem{remark}[theorem]{Remark}

\title{One Obstacle, Two Unsolved Problems:\\
How Collatz and Navier--Stokes Point to the Same Missing Language}

\author{Pablo Nogueira Grossi\\
G6 LLC · Newark, NJ · 2026\\
\texttt{pablogrossi@hotmail.com}\\
ORCID: 0009-0000-6496-2186}

\date{April 2026 (preprint)}

\begin{document}

\maketitle

\begin{abstract}
We propose that the Collatz conjecture is a canonical instance of a
\emph{dm\textsuperscript{3}-system}: a dynamical system governed by the
operator algebra $\mathcal{G} = U \circ F \circ K \circ C$ (Unfold,
Fold, Curvature, Compression) acting on a contact manifold.
Under this framing the mean-contraction property of the Collatz map
with coefficient $c = 3$ is not a numerical accident but a structural
fingerprint of the operator grammar — the same fingerprint that selects
wavenumber $n = 6$ in Saturn's polar hexagon, governs
reconnection-onset in plasma instabilities, and drives fold transitions
in biological morphogenesis.

We do not claim a proof of the conjecture.
We claim \emph{structural visibility prior to axiomatisation}: the
dm\textsuperscript{3} framework makes the relevant invariants explicit,
identifies the precise gap (Bridge~0, the local-to-global dissipation
inequality) that separates observable contraction from a closure theorem,
and places Collatz in the same universality class as the Navier--Stokes
regularity problem.
We invite fluid dynamicists, harmonic analysts, number theorists, and
applied mathematicians to supply the local hypotheses and empirical
diagnostics needed to formalise the missing language.
\end{abstract}

\tableofcontents
\bigskip

%% ------------------------------------------------------------------
\section{Introduction}
%% ------------------------------------------------------------------

The Collatz map $T : \mathbb{N} \to \mathbb{N}$,
\[
  T(n) = \begin{cases} n/2 & \text{if } n \text{ is even,} \\
                        3n+1 & \text{if } n \text{ is odd,}
         \end{cases}
\]
has resisted proof for nearly a century despite the empirical
universality of the cycle $\{4 \to 2 \to 1\}$.
Standard analytic approaches — density arguments, $p$-adic methods,
ergodic theory — each establish partial results but fall short of a
global closure theorem.

We argue that this persistent gap is not merely a technical obstacle
but a symptom of a \emph{missing formal language}: a higher-order logic
that can translate local contraction estimates into global convergence
statements across both discrete arithmetic and continuous analysis.
The same gap appears in the Navier--Stokes regularity problem, and
framing both as instances of a single obstacle clarifies why parallel
progress across domains is not only plausible but structurally expected.

The paper proceeds as follows.
Section~2 recalls the dm\textsuperscript{3} framework and the operator
grammar $\mathcal{G} = U \circ F \circ K \circ C$.
Section~3 presents Saturn's polar vortex as an empirical certificate
for the framework.
The new Section~3.1 identifies the common structural obstacle shared by
Collatz and Navier--Stokes.
Section~4 analyses the $c = 3$ coefficient as a triad fingerprint.
Sections~5--9 develop the Collatz dm\textsuperscript{3} identification,
the open axioms (Bridge~0), the Lean~4 formalisation programme, and
the cross-disciplinary invitation.

%% ------------------------------------------------------------------
\section{The dm\textsuperscript{3} Framework}
%% ------------------------------------------------------------------

\subsection{The Operator Algebra}

\begin{definition}[Generative Operator]
Let $(X, g)$ be a compact Riemannian manifold with contact form
$\alpha$.  The \emph{generative operator} is the composition
\[
  \mathcal{G} = U \circ F \circ K \circ C : X \to X,
\]
where
\begin{enumerate}[label=(\roman*)]
  \item $C$ (\emph{Compression}): a Lipschitz map onto a
        lower-dimensional submanifold, $\kappa = \mathrm{Lip}(C) < 1$;
  \item $K$ (\emph{Curvature}): a Riemannian-curvature flow driving
        orbits toward the critical threshold $\kappa^*$;
  \item $F$ (\emph{Fold}): a Whitney $A_1$ singularity triggered at
        $\|\kappa\| = \kappa^*$, where $\det(Df) \to 0$;
  \item $U$ (\emph{Unfold}): gradient descent on a Lyapunov functional
        $\Phi$, selecting the stable attractor ($\nabla\Phi = 0$).
\end{enumerate}
\end{definition}

\subsection{The Canonical dm\textsuperscript{3} Invariants}

The toy model on $M = S \times \mathbb{R}$ with contact form
$\alpha = dz - r^2\,d\theta$ realises the canonical triple
\[
  \iota(\mathfrak{D}) = (T^*, \mu_{\max}, \tau) = (2\pi,\,-2,\,2).
\]
The structural stability radius is $\varepsilon_0 = 1/3$, and the
noise-tolerance product satisfies $\tau \cdot \varepsilon_0 = 2/3$.

\begin{remark}[Epistemic status]
Axioms~1--6 of the dm\textsuperscript{3} framework are structural
definitions.  Axioms~7--8 (no divergent orbits; unique attractor) are
open for the Collatz map — they \emph{are} the conjecture.
All claims in this paper are conditional on those open axioms unless
explicitly stated otherwise.
\end{remark}

%% ------------------------------------------------------------------
\section{The Polar Vortex as Empirical Certificate}
%% ------------------------------------------------------------------

Saturn's north polar hexagon constitutes the strongest available
empirical certificate for the dm\textsuperscript{3} framework.
The mapping is clean and over-determined:

\begin{enumerate}
  \item \textbf{C (Compression):} The 3-D atmosphere compresses into a
        quasi-2D zonal jet layer.
  \item \textbf{K (Curvature):} Rossby wave curvature accumulates,
        selecting wavenumber $n = 6$ — not 5, not 7 — via the
        constraint $\kappa^* = \varepsilon_0 = 1/3$.
  \item \textbf{F (Fold):} The Whitney fold locks six sharp corners
        with sub-kilometer precision.
  \item \textbf{U (Unfold):} Gradient descent on $\Phi$ yields the
        persistent $40^+$-year fixed point observed from Voyager through
        Cassini (1980--2020).
\end{enumerate}

The $c = 3$ triad fingerprint (Section~4) reappears here: the
Rossby-wave triad resonance selects precisely $n = 6 = 2 \times 3$
and suppresses $n = 4, 5, 7, 8$.  The polar vortex is therefore both a
test case and a parameter-estimation laboratory for the discrete
arithmetic application.

%% ------------------------------------------------------------------
\subsection{The Same Gap Across Two Famous Unsolved Problems}
%% ------------------------------------------------------------------

Two of the longest-standing open problems in mathematics --- the
Collatz conjecture in discrete arithmetic and the Navier--Stokes
regularity problem in fluid dynamics --- share a common structural
obstacle we call \textbf{Bridge~0}.  In both settings a local mechanism
appears to dissipate or contract at small scales, yet the available
analytic language fails to convert that local control into a robust
global closure theorem.  For Navier--Stokes the obstruction is endpoint
control of nonlinear fluxes across scales; for Collatz it is the lack
of a discrete analogue of the local-to-global dissipation inequality
(the mean-contraction bridge).  Framing these as instances of the same
missing step clarifies why progress in one domain can illuminate the
other.

A higher-order logic that unifies the operator grammar
$G = U \circ F \circ K \circ C$ across continuous manifolds and
discrete arithmetic would turn both problems into corollaries of a
single closure principle.  Such a formalism would supply the discrete
analogues of Lyapunov functionals, contact forms, and scale-invariant
closure theorems needed to convert local contraction estimates into
global convergence or regularity.  This is not a claim that the
problems are trivial; it is a claim that the \emph{same} missing formal
ingredient blocks both, and that constructing that ingredient would
unlock parallel advances across domains.

This is therefore a cross-disciplinary programme.  Fluid dynamicists,
harmonic analysts, number theorists, plasma physicists, and applied
mathematicians each encounter manifestations of the same generative
operator grammar at their scale.  We invite domain experts to
(a)~identify the precise local control statements they already use in
practice, (b)~translate those statements into candidate axioms for a
higher-order language, and (c)~test the axioms against empirical
diagnostics (Fourier decay, operator numerics, exceptional sets).
Your expertise is the missing ingredient.

\begin{quote}
\emph{This programme is conditional and collaborative --- we do not
claim a proof.  We invite experts from each domain to contribute
precise local hypotheses and empirical diagnostics so the higher-order
language can be formalised and tested.}
\end{quote}

%% ------------------------------------------------------------------
\section{The \texorpdfstring{$c = 3$}{c=3} Coefficient as Triad Fingerprint}
%% ------------------------------------------------------------------

\subsection{Mean-Contraction Analysis}

For the generalised Collatz map with odd-branch coefficient $c$,
\[
  T_c(n) = \begin{cases} n/2 & n \text{ even,} \\ cn+1 & n \text{ odd,} \end{cases}
\]
the mean step ratio over one even--odd pair is
\[
  \rho(c) = \frac{c}{2 \cdot 2} = \frac{c}{4}.
\]
Mean contraction requires $\rho(c) < 1$, i.e.\ $c < 4$.
Since $c$ must be odd and $c \geq 1$:
\[
  c = 3 \quad \Longrightarrow \quad \rho(3) = \tfrac{3}{4} < 1 \quad
  \text{(contraction);}
\]
\[
  c = 5 \quad \Longrightarrow \quad \rho(5) = \tfrac{5}{4} > 1 \quad
  \text{(divergence on average).}
\]

\begin{proposition}[Minimality of $c = 3$]
$c = 3$ is the unique minimal odd positive integer producing mean
contraction in $T_c$.
\end{proposition}

\subsection{Triad Resonance}

In the operator-grammar language, $c = 3$ encodes the minimum
three-element resonant triad (even $\to$ odd $\to$ attractor) that
closes under $\mathcal{G}$.  The same triad structure selects $n = 6$
in the Saturn hexagon ($6 = 2 \times 3$) and the $4:1$ Arnold tongue
in Schumann resonance ($f_4 \approx 33.5$ Hz $\sim 2 \times g^6$
with $g^6 = 33$).

%% ------------------------------------------------------------------
\section{Collatz as a dm\textsuperscript{3}-System}
%% ------------------------------------------------------------------

\subsection{Operator Identification}

We identify the Collatz dynamics with the dm\textsuperscript{3}
operator sequence as follows:

\begin{center}
\begin{tabular}{lll}
\hline
\textbf{Operator} & \textbf{Collatz instantiation} & \textbf{Formal role} \\
\hline
$C$ & Even-step halving: $n \mapsto n/2$ & Lipschitz compression $\kappa = 1/2$ \\
$K$ & Parity curvature: odd/even alternation & Threshold accumulation \\
$F$ & Odd-branch expansion: $n \mapsto 3n+1$ & Whitney $A_1$ local model \\
$U$ & Descent to $\{4,2,1\}$ attractor & Gradient descent on $\Phi$ \\
\hline
\end{tabular}
\end{center}

\subsection{Structural Hypothesis}

\begin{conjecture}[Collatz dm\textsuperscript{3} identification]
The Collatz dynamical system $(\mathbb{N}, T)$ is a dm\textsuperscript{3}
system with canonical invariants
\[
  (T^*, \mu_{\max}, \tau) = (2\pi,\,-2,\,2), \quad \varepsilon_0 = 1/3,
\]
where $T^*$ is the attractor period (the $\{4 \to 2 \to 1\}$ cycle),
$\mu_{\max}$ is the maximal Lyapunov exponent of the discrete orbit,
and $\tau = \sqrt{c/\kappa} = \sqrt{3/(1/2)} \cdot (2/3) = 2$ is the
embodiment threshold.
\end{conjecture}

This conjecture does not imply a proof of convergence.  It identifies
the precise structural parameters that a proof would need to verify.

%% ------------------------------------------------------------------
\section{Bridge~0: The Open Gap}
%% ------------------------------------------------------------------

\subsection{Statement of Bridge~0}

\begin{definition}[Bridge~0]
\emph{Bridge~0} is the local-to-global dissipation inequality:
a formal statement of the form
\[
  \text{``local contraction } \Rightarrow \text{ global convergence''}
\]
valid in the discrete arithmetic setting of $T$.
\end{definition}

More precisely, the mean-contraction property $\rho(3) = 3/4 < 1$
(established by direct computation) must be upgraded to an inequality
of the form
\[
  \Phi(T^k(n)) \leq (1 - \delta)^k \Phi(n) \quad \forall\, n \in \mathbb{N},
\]
for some Lyapunov functional $\Phi : \mathbb{N} \to \mathbb{R}_{\geq 0}$
and $\delta > 0$ uniform in $n$.

\subsection{Why Standard Tools Fail}

\begin{itemize}
  \item \textbf{Density / ergodic arguments} establish almost-sure
        convergence but cannot rule out a single exceptional orbit.
  \item \textbf{$p$-adic methods} control the $2$-adic and $3$-adic
        structure locally but lack a bridge from local to global
        $\ell^\infty$ control.
  \item \textbf{Harmonic analysis} (Fourier decay, Tao's logarithmic
        density result) shows that almost all orbits eventually reach
        below any bound, but the endpoint $\ell^\infty$ estimate
        remains open.
\end{itemize}

Each of these corresponds, in continuous analysis, to the endpoint
regularity problem for Navier--Stokes: local $L^2$ energy dissipation
is controlled, but the $L^\infty$ or $\dot{H}^{1/2}$ endpoint is not.
This structural parallel is Bridge~0's primary diagnostic content.

%% ------------------------------------------------------------------
\section{The Lean~4 Formalisation Programme (AXLE)}
%% ------------------------------------------------------------------

The AXLE repository (\url{https://github.com/TOTOGT/PabloNogueira-dm3-lab})
contains Lean~4 formalisations of the structural components that \emph{do}
admit formal proof:

\begin{enumerate}
  \item \texttt{TOGT.noiseTolerance} — $\tau \cdot \varepsilon_0 = 2/3$
  \item \texttt{TOGT.stabilityRadius\_eq} — $\varepsilon_0 = 1/3$
  \item \texttt{TOGT.crystal\_aspect\_ratio} — height/base $= 66$
  \item \texttt{TOGT.aspect\_ratio\_encodes\_invariants} — $66 = g^6 \cdot \tau$
  \item \texttt{TOGT.regeneration\_unbounded} — hierarchy is unbounded
\end{enumerate}

Axioms~7--8 (no divergent orbits; unique attractor) are designated
\texttt{AXLE Target 5} and remain open, consistent with the Collatz
conjecture itself.

%% ------------------------------------------------------------------
\section{Empirical Diagnostics}
%% ------------------------------------------------------------------

\subsection{Fourier Decay of Exceptional Sets}

Following Tao~\cite{Tao2022}, define the exceptional set
$E_N = \{n \leq N : T^k(n) > n \text{ for all } k \leq K(n)\}$.
The dm\textsuperscript{3} framework predicts:
\[
  |E_N| = O(N^{1 - \delta}) \quad \text{for some } \delta > 0.
\]
This is consistent with Tao's result $|E_N| = O(N / \log N)$ but does
not sharpen it without Bridge~0.

\subsection{Operator Numerics}

Python simulations in \texttt{simulations/simple\_to\_operator.py}
compute the dm\textsuperscript{3} metric
$\tau = \sqrt{c/\kappa}$ and Lyapunov exponents for orbit segments of
length up to $10^6$.  Preliminary results confirm $\rho(3) = 3/4$
and detect no divergent segment up to $n = 2^{68}$ (Oliveira--Leal
census).

%% ------------------------------------------------------------------
\section{The Cross-Disciplinary Invitation}
%% ------------------------------------------------------------------

The programme described here is explicitly collaborative and conditional.
We identify a structural obstacle (Bridge~0) and a formal language that
makes the obstacle explicit.  We do not claim to have resolved either
the Collatz conjecture or the Navier--Stokes regularity problem.

We invite:

\begin{itemize}
  \item \textbf{Number theorists and harmonic analysts} to identify
        candidate Lyapunov functionals $\Phi$ and test whether they
        satisfy the discrete dissipation inequality.
  \item \textbf{Fluid dynamicists} to translate their local energy
        estimates into axiom candidates for the higher-order language,
        and to test those axioms against Navier--Stokes numerics.
  \item \textbf{Plasma physicists} to contribute reconnection-rate
        estimates as empirical diagnostics for the fold-operator
        identification.
  \item \textbf{Applied mathematicians and computer scientists} to
        extend the Lean~4 formalisation (AXLE) to cover new axioms as
        they are proposed.
\end{itemize}

The repository is MIT licensed.  Domain-mapping contributions follow
the template in \texttt{mappings/domain\_mappings.md}.

%% ------------------------------------------------------------------
\section*{Acknowledgements}
%% ------------------------------------------------------------------

The dm\textsuperscript{3} framework was developed in the \emph{Principia
Orthogona} series (Volumes I--VI, G6 LLC).  Saturn hexagon parameter
data are drawn from Cassini VIMS and ISS archives (NASA/ESA/ASI).
Lean~4 formalisation uses Mathlib4.  The connectome simulation uses the
FlyWire adult female fly connectome (\url{https://codex.flywire.ai}).

%% ------------------------------------------------------------------
\begin{thebibliography}{99}

\bibitem{Collatz1937}
L.~Collatz,
\emph{Unpublished problem presented at the 7th International Congress of
Mathematicians}, Cambridge, 1950.

\bibitem{Tao2022}
T.~Tao,
\emph{Almost all orbits of the Collatz map attain almost bounded values},
Forum of Mathematics, Pi, \textbf{10} (2022), e12.
\url{https://doi.org/10.1017/fmp.2022.8}

\bibitem{Fefferman2000}
C.~L.~Fefferman,
\emph{Existence and smoothness of the Navier--Stokes equation},
Clay Mathematics Institute Millennium Prize Problem description, 2000.

\bibitem{Lagarias2010}
J.~C.~Lagarias (ed.),
\emph{The Ultimate Challenge: The 3x+1 Problem},
American Mathematical Society, Providence, RI, 2010.

\bibitem{PrincipalOrthogona}
P.~Nogueira Grossi,
\emph{Principia Orthogona}, Volumes I--VI,
G6 LLC, Newark, NJ, 2024--2026.

\bibitem{Zenodo2026}
P.~Nogueira Grossi,
\emph{The Collatz Conjecture as a Canonical dm\textsuperscript{3}-System:
A Structural Framework for Decidability},
Zenodo, 2026.
\url{https://doi.org/10.5281/zenodo.19117400}

\end{thebibliography}

\end{document}
